\documentclass[10pt]{extarticle}
\usepackage[a4paper,margin=1.5cm]{geometry}
\usepackage[T1]{fontenc}
\usepackage{amsmath,amssymb}
\usepackage{graphicx}
\usepackage[export]{adjustbox}
\usepackage{multicol}
\usepackage{enumitem}
\usepackage{booktabs,array,tabularx}
\setlength{\columnsep}{0.7cm}
\setlength{\parindent}{0pt}
\setlength{\parskip}{3pt}
\pagestyle{empty}
\makeatletter
% figure registry (figdefs.tex is generated by the extract step)
\newcommand{\qfigdef}[4]{\@namedef{qf@file@#1}{#2}\@namedef{qf@w@#1}{#3}\@namedef{qf@wide@#1}{#4}}
\newcommand{\qf@check}[1]{\@ifundefined{qf@file@#1}{\PackageError{score2}{Unknown figure id #1}{}}{}}
\newcommand{\qf@img}[4][]{\includegraphics[width=#2\dimexpr\@nameuse{qf@w@#4}pt\relax,max width=#3,max height=0.45\textheight,#1]{\@nameuse{qf@file@#4}}}
% Full-width questions: finalize wraps every question that contains a wide figure in \qfullbegin...\qfullend,
% so the WHOLE question (stem, figures, options) is set across the page in reading order, then the two columns
% resume. The question's head - its number, stem and first figure (\qfullsplit marks the end of the head) -
% is kept on one page, so a number is never left at a page foot with its figure overleaf; the rest may break
% across pages normally, so no big gaps are left. Without a \qfullsplit, the whole question is kept together
% if it is under 45% of a page. Never write these yourself.
\newif\ifqf@full
\newif\ifqf@boxed
\newsavebox{\qf@box}
\newcommand{\qf@open}{\global\qf@fulltrue\par\global\qf@boxedtrue
  \setbox\qf@box\vbox\bgroup\hsize\textwidth\linewidth\textwidth}
\newcommand{\qf@emit}[1]{\par\egroup\global\qf@boxedfalse
  \ifdim\dimexpr\ht\qf@box+\dp\qf@box\relax<#1\textheight\noindent\usebox{\qf@box}\par\else\unvbox\qf@box\fi}
\newcommand{\qfullsplit}{\ifqf@boxed\qf@emit{0.85}\fi}
\newcommand{\qf@close}{\ifqf@boxed\qf@emit{0.45}\fi\global\qf@fullfalse}
\newcommand{\qfullbegin}{\end{multicols}\qf@open}
\newcommand{\qfullend}{\qf@close\begin{multicols}{2}}
\newcommand{\qfullstart}{\qf@open}
\newcommand{\qfulldone}{\qf@close}
% \qfig[scale]{id}: block figure, centred, directly where it is written (size set by the tool).
% Fallback if a wide figure is not inside a full-width question: break out of the columns in place - no floats.
\newcommand{\qfig}[2][1]{\qf@check{#2}%
  \ifqf@full\qf@blk{#1}{#2}\else
  \ifnum\@nameuse{qf@wide@#2}=1 \expandafter\@firstoftwo\else\expandafter\@secondoftwo\fi
  {\end{multicols}\begin{center}\qf@img{#1}{\textwidth}{#2}\end{center}\begin{multicols}{2}}%
  {\qf@blk{#1}{#2}}\fi}
\newcommand{\qf@blk}[2]{\par\smallskip{\centering\qf@img{#1}{\linewidth}{#2}\par}\smallskip}
% \qfigi[scale]{id}: inline figure (inside option lists / tables / minipages); never breaks columns.
\newcommand{\qfigi}[2][1]{\qf@check{#2}\qf@img[valign=c]{#1}{\linewidth}{#2}}
\makeatother
% structure macros - lint reads these to check completeness and figure placement
\newcommand{\papersec}[1]{\par\bigskip{\centering\Large\bfseries PAPER-#1\par}\medskip}
\newcommand{\partsec}[2]{\par\medskip\noindent{\large\bfseries PART-#1: #2}\par\smallskip}
\newcommand{\sect}[1]{\par\smallskip\noindent\textbf{#1}\par}
\newcommand{\qn}[1]{\par\smallskip\noindent\textbf{#1.}~}
\newcommand{\paraq}[2]{\par\smallskip\noindent\textbf{Paragraph for Questions #1 and #2}\par}
\newenvironment{qoptions}{\begin{itemize}[leftmargin=1.8em,itemsep=1pt,topsep=2pt,label={}]}{\end{itemize}}
\newcommand{\opt}[2]{\item[(#1)] #2}
% figure registry (generated by score2_tools.py finalize)
\qfigdef{p1_phy_s1_q01_1}{TOAS - Test 13 (24 Jul 2022)-img/p1_phy_s1_q01_1.png}{31.8}{0}
\qfigdef{p1_phy_s1_q02_1}{TOAS - Test 13 (24 Jul 2022)-img/p1_phy_s1_q02_1.png}{100.4}{0}
\qfigdef{p1_phy_s1_q05_1}{TOAS - Test 13 (24 Jul 2022)-img/p1_phy_s1_q05_1.png}{62.0}{0}
\qfigdef{p1_phy_s1_q06_1}{TOAS - Test 13 (24 Jul 2022)-img/p1_phy_s1_q06_1.png}{111.1}{0}
\qfigdef{p1_phy_s2_para7_8_1}{TOAS - Test 13 (24 Jul 2022)-img/p1_phy_s2_para7_8_1.png}{99.7}{0}
\qfigdef{p1_phy_s2_para9_10_1}{TOAS - Test 13 (24 Jul 2022)-img/p1_phy_s2_para9_10_1.png}{109.4}{0}
\qfigdef{p1_phy_s3_q01_1}{TOAS - Test 13 (24 Jul 2022)-img/p1_phy_s3_q01_1.png}{111.2}{0}
\qfigdef{p1_phy_s3_q02_1}{TOAS - Test 13 (24 Jul 2022)-img/p1_phy_s3_q02_1.png}{89.9}{0}
\qfigdef{p1_phy_s3_q03_1}{TOAS - Test 13 (24 Jul 2022)-img/p1_phy_s3_q03_1.png}{113.1}{0}
\qfigdef{p1_phy_s3_q04_1}{TOAS - Test 13 (24 Jul 2022)-img/p1_phy_s3_q04_1.png}{116.3}{0}
\qfigdef{p1_phy_s3_q05_1}{TOAS - Test 13 (24 Jul 2022)-img/p1_phy_s3_q05_1.png}{76.8}{0}
\qfigdef{p1_phy_s3_q06_1}{TOAS - Test 13 (24 Jul 2022)-img/p1_phy_s3_q06_1.png}{87.4}{0}
\qfigdef{p1_phy_s3_q07_1}{TOAS - Test 13 (24 Jul 2022)-img/p1_phy_s3_q07_1.png}{122.2}{0}
\qfigdef{p1_phy_s3_q08_1}{TOAS - Test 13 (24 Jul 2022)-img/p1_phy_s3_q08_1.png}{117.6}{0}
\qfigdef{p1_chem_s1_q03_A}{TOAS - Test 13 (24 Jul 2022)-img/p1_chem_s1_q03_A.png}{158.0}{0}
\qfigdef{p1_chem_s1_q03_B}{TOAS - Test 13 (24 Jul 2022)-img/p1_chem_s1_q03_B.png}{126.4}{0}
\qfigdef{p1_chem_s1_q03_C}{TOAS - Test 13 (24 Jul 2022)-img/p1_chem_s1_q03_C.png}{124.1}{0}
\qfigdef{p1_chem_s1_q03_D}{TOAS - Test 13 (24 Jul 2022)-img/p1_chem_s1_q03_D.png}{127.1}{0}
\qfigdef{p1_chem_s1_q04_B}{TOAS - Test 13 (24 Jul 2022)-img/p1_chem_s1_q04_B.png}{183.5}{0}
\qfigdef{p1_chem_s1_q04_3}{TOAS - Test 13 (24 Jul 2022)-img/p1_chem_s1_q04_3.png}{97.9}{0}
\qfigdef{p1_chem_s2_q10_A}{TOAS - Test 13 (24 Jul 2022)-img/p1_chem_s2_q10_A.png}{97.7}{0}
\qfigdef{p1_chem_s2_q10_B}{TOAS - Test 13 (24 Jul 2022)-img/p1_chem_s2_q10_B.png}{100.4}{0}
\qfigdef{p1_chem_s2_q10_C}{TOAS - Test 13 (24 Jul 2022)-img/p1_chem_s2_q10_C.png}{113.6}{0}
\qfigdef{p1_chem_s2_q10_D}{TOAS - Test 13 (24 Jul 2022)-img/p1_chem_s2_q10_D.png}{108.3}{0}
\qfigdef{p1_chem_s3_q03_1}{TOAS - Test 13 (24 Jul 2022)-img/p1_chem_s3_q03_1.png}{210.9}{0}
\qfigdef{p1_chem_s3_q05_1}{TOAS - Test 13 (24 Jul 2022)-img/p1_chem_s3_q05_1.png}{258.7}{1}
\qfigdef{p2_phy_s1_q01_1}{TOAS - Test 13 (24 Jul 2022)-img/p2_phy_s1_q01_1.png}{133.0}{0}
\qfigdef{p2_phy_s1_q03_1}{TOAS - Test 13 (24 Jul 2022)-img/p2_phy_s1_q03_1.png}{94.3}{0}
\qfigdef{p2_phy_s1_q07_1}{TOAS - Test 13 (24 Jul 2022)-img/p2_phy_s1_q07_1.png}{113.7}{0}
\qfigdef{p2_phy_s1_q08_1}{TOAS - Test 13 (24 Jul 2022)-img/p2_phy_s1_q08_1.png}{121.7}{0}
\qfigdef{p2_phy_s1_q09_1}{TOAS - Test 13 (24 Jul 2022)-img/p2_phy_s1_q09_1.png}{63.9}{0}
\qfigdef{p2_phy_s1_q12_1}{TOAS - Test 13 (24 Jul 2022)-img/p2_phy_s1_q12_1.png}{94.2}{0}
\qfigdef{p2_phy_s1_q15_1}{TOAS - Test 13 (24 Jul 2022)-img/p2_phy_s1_q15_1.png}{152.8}{0}
\qfigdef{p2_phy_s1_q16_1}{TOAS - Test 13 (24 Jul 2022)-img/p2_phy_s1_q16_1.png}{95.4}{0}
\qfigdef{p2_phy_s2_q01_1}{TOAS - Test 13 (24 Jul 2022)-img/p2_phy_s2_q01_1.png}{137.6}{0}
\qfigdef{p2_phy_s2_q02_1}{TOAS - Test 13 (24 Jul 2022)-img/p2_phy_s2_q02_1.png}{132.5}{0}
\qfigdef{p2_phy_s2_q04_1}{TOAS - Test 13 (24 Jul 2022)-img/p2_phy_s2_q04_1.png}{105.4}{0}
\qfigdef{p2_phy_s2_q05_1}{TOAS - Test 13 (24 Jul 2022)-img/p2_phy_s2_q05_1.png}{148.7}{0}
\qfigdef{p2_phy_s2_q07_1}{TOAS - Test 13 (24 Jul 2022)-img/p2_phy_s2_q07_1.png}{149.1}{0}
\qfigdef{p2_phy_s2_q10_1}{TOAS - Test 13 (24 Jul 2022)-img/p2_phy_s2_q10_1.png}{130.8}{0}
\qfigdef{p2_chem_s1_q03_A}{TOAS - Test 13 (24 Jul 2022)-img/p2_chem_s1_q03_A.png}{63.2}{0}
\qfigdef{p2_chem_s1_q03_B}{TOAS - Test 13 (24 Jul 2022)-img/p2_chem_s1_q03_B.png}{63.4}{0}
\qfigdef{p2_chem_s1_q03_C}{TOAS - Test 13 (24 Jul 2022)-img/p2_chem_s1_q03_C.png}{63.4}{0}
\qfigdef{p2_chem_s1_q03_D}{TOAS - Test 13 (24 Jul 2022)-img/p2_chem_s1_q03_D.png}{63.4}{0}
\qfigdef{p2_chem_s1_q04_1}{TOAS - Test 13 (24 Jul 2022)-img/p2_chem_s1_q04_1.png}{100.1}{0}
\qfigdef{p2_chem_s1_q07_1}{TOAS - Test 13 (24 Jul 2022)-img/p2_chem_s1_q07_1.png}{132.7}{0}
\qfigdef{p2_chem_s1_q07_2}{TOAS - Test 13 (24 Jul 2022)-img/p2_chem_s1_q07_2.png}{132.2}{0}
\qfigdef{p2_chem_s1_q08_A}{TOAS - Test 13 (24 Jul 2022)-img/p2_chem_s1_q08_A.png}{126.9}{0}
\qfigdef{p2_chem_s1_q08_B}{TOAS - Test 13 (24 Jul 2022)-img/p2_chem_s1_q08_B.png}{134.7}{0}
\qfigdef{p2_chem_s1_q08_C}{TOAS - Test 13 (24 Jul 2022)-img/p2_chem_s1_q08_C.png}{134.7}{0}
\qfigdef{p2_chem_s1_q08_D}{TOAS - Test 13 (24 Jul 2022)-img/p2_chem_s1_q08_D.png}{138.9}{0}
\qfigdef{p2_chem_s1_q09_1}{TOAS - Test 13 (24 Jul 2022)-img/p2_chem_s1_q09_1.png}{130.8}{0}
\qfigdef{p2_chem_s1_q09_A}{TOAS - Test 13 (24 Jul 2022)-img/p2_chem_s1_q09_A.png}{64.6}{0}
\qfigdef{p2_chem_s1_q09_B}{TOAS - Test 13 (24 Jul 2022)-img/p2_chem_s1_q09_B.png}{43.9}{0}
\qfigdef{p2_chem_s1_q09_C}{TOAS - Test 13 (24 Jul 2022)-img/p2_chem_s1_q09_C.png}{38.0}{0}
\qfigdef{p2_chem_s1_q09_D}{TOAS - Test 13 (24 Jul 2022)-img/p2_chem_s1_q09_D.png}{43.9}{0}
\qfigdef{p2_chem_s1_q10_1}{TOAS - Test 13 (24 Jul 2022)-img/p2_chem_s1_q10_1.png}{153.3}{0}
\qfigdef{p2_chem_s1_q10_A}{TOAS - Test 13 (24 Jul 2022)-img/p2_chem_s1_q10_A.png}{133.9}{0}
\qfigdef{p2_chem_s1_q10_B}{TOAS - Test 13 (24 Jul 2022)-img/p2_chem_s1_q10_B.png}{135.4}{0}
\qfigdef{p2_chem_s1_q10_C}{TOAS - Test 13 (24 Jul 2022)-img/p2_chem_s1_q10_C.png}{134.9}{0}
\qfigdef{p2_chem_s1_q10_D}{TOAS - Test 13 (24 Jul 2022)-img/p2_chem_s1_q10_D.png}{114.6}{0}
\qfigdef{p2_chem_s1_q11_1}{TOAS - Test 13 (24 Jul 2022)-img/p2_chem_s1_q11_1.png}{149.1}{0}
\qfigdef{p2_chem_s1_q11_A}{TOAS - Test 13 (24 Jul 2022)-img/p2_chem_s1_q11_A.png}{49.9}{0}
\qfigdef{p2_chem_s1_q11_B}{TOAS - Test 13 (24 Jul 2022)-img/p2_chem_s1_q11_B.png}{60.0}{0}
\qfigdef{p2_chem_s1_q11_C}{TOAS - Test 13 (24 Jul 2022)-img/p2_chem_s1_q11_C.png}{52.4}{0}
\qfigdef{p2_chem_s1_q11_D}{TOAS - Test 13 (24 Jul 2022)-img/p2_chem_s1_q11_D.png}{58.3}{0}
\qfigdef{p2_chem_s1_q12_1}{TOAS - Test 13 (24 Jul 2022)-img/p2_chem_s1_q12_1.png}{141.6}{0}
\qfigdef{p2_chem_s1_q12_A}{TOAS - Test 13 (24 Jul 2022)-img/p2_chem_s1_q12_A.png}{24.0}{0}
\qfigdef{p2_chem_s1_q12_B}{TOAS - Test 13 (24 Jul 2022)-img/p2_chem_s1_q12_B.png}{59.3}{0}
\qfigdef{p2_chem_s1_q12_C}{TOAS - Test 13 (24 Jul 2022)-img/p2_chem_s1_q12_C.png}{26.5}{0}
\qfigdef{p2_chem_s1_q12_D}{TOAS - Test 13 (24 Jul 2022)-img/p2_chem_s1_q12_D.png}{59.3}{0}
\qfigdef{p2_chem_s1_q17_1}{TOAS - Test 13 (24 Jul 2022)-img/p2_chem_s1_q17_1.png}{154.5}{0}
\qfigdef{p2_chem_s2_q01_1}{TOAS - Test 13 (24 Jul 2022)-img/p2_chem_s2_q01_1.png}{82.0}{0}
\qfigdef{p2_chem_s2_q04_1}{TOAS - Test 13 (24 Jul 2022)-img/p2_chem_s2_q04_1.png}{99.9}{0}
\qfigdef{p2_chem_s2_q05_1}{TOAS - Test 13 (24 Jul 2022)-img/p2_chem_s2_q05_1.png}{53.2}{0}
\qfigdef{p2_chem_s2_q10_1}{TOAS - Test 13 (24 Jul 2022)-img/p2_chem_s2_q10_1.png}{123.7}{0}
\raggedcolumns
\makeatletter\@beginparpenalty=10000\makeatother
\newcolumntype{L}[1]{>{\raggedright\arraybackslash}p{#1}}
\begin{document}
\begin{center}{\Large\bfseries TOAS --- Test 13 (24 Jul 2022)}\end{center}
\begin{multicols}{2}
\raggedcolumns
\papersec{1}
\partsec{1}{PHYSICS}
\sect{SECTION-I(i)}
\qn{1} A mass $m$ is hung on an ideal massless spring. Another equal mass is connected to the other end of the spring. The whole system is at rest. At $t=0$, $m$ is released and the system falls freely under gravity. Assume that natural length of the spring is $L_0$, its initial stretched length is $L$ and the acceleration due to gravity is $g$. What is distance between masses as function of time.
\qfig{p1_phy_s1_q01_1}
\begin{qoptions}
\opt{A}{$L_0+(L-L_0)\cos\sqrt{\dfrac{2K}{m}}\,t$}
\opt{B}{$L_0+(L-L_0)\cos\sqrt{\dfrac{K}{m}}\,t$}
\opt{C}{$L_0-2(L+2L_0)\cos\sqrt{\dfrac{2K}{m}}\,t$}
\opt{D}{$L_0+(L-L_0)\sin\sqrt{\dfrac{2K}{m}}\,t$}
\end{qoptions}
\qn{2} Consider a situation in which a uniform gravitation field is present below line AB and no gravity above line AB. Then the speed with which a rod shown in the figure will strike earth surface when released from the position shown is
\qfig{p1_phy_s1_q02_1}
\begin{qoptions}
\opt{A}{$\sqrt{g\ell}$}
\opt{B}{$2\sqrt{g\ell}$}
\opt{C}{$3\sqrt{g\ell}$}
\opt{D}{$4\sqrt{g\ell}$}
\end{qoptions}
\qn{3} A non-conducting solid sphere of radius 30cm and relative permittivity 1 has the volume charge density $\rho=\left(\dfrac{5}{2\pi}\,\mu\text{C/m}^3\right)\left(1-\dfrac{r}{30\text{cm}}\right)$ where $r$ (in cm) is the distance of any point in the space from the centre of the sphere. Then choose the correct statement(s).
\begin{qoptions}
\opt{A}{Density of electric field lines is maximum at $r=25$ cm.}
\opt{B}{Density of electric field lines is maximum at $r=20$ cm.}
\opt{C}{Electrical energy stored in the space $r=30$ cm to $r\to\infty$ is same as that of a solid non-conducting sphere carrying charges of 0.0225 $\mu$C and radii 30 cm having uniform volume charge density in the same space $r=30$ cm to $r\to\infty$.}
\opt{D}{Variation of magnitude of electric field due to the sphere with respect to radial distance $r$ is rectangular hyperbola in the region $r=30$ cm to $r\to\infty$.}
\end{qoptions}
\qn{4} A closed vessel contains a mixture of two diatomic gases A and B. Molar mass of A is 16 times that of B and mass of gas A, contained in the vessel is 2 times that of B :-
\begin{qoptions}
\opt{A}{Pressure exerted by gas B is eight times that of exerted by gas A.}
\opt{B}{Number of molecules of gas B in the cylinder is eight times that of gas A.}
\opt{C}{Average kinetic energy per molecule of gas A is equal to that of gas B.}
\opt{D}{Root mean square value of translational speed of gas B is four times that of A.}
\end{qoptions}
\qn{5} In the cyclic process shown in the given figure, $\Delta U_1$ and $\Delta U_2$ represent the change in internal energy in process A and B, respectively. If $\Delta Q$ be the net heat given to the system in the process and $\Delta W$ be the work done by the system in the process, then
\qfig{p1_phy_s1_q05_1}
\begin{qoptions}
\opt{A}{$\Delta U_1+\Delta U_2=0$}
\opt{B}{$\Delta U_1-\Delta U_2=0$}
\opt{C}{$\Delta Q-\Delta W=0$}
\opt{D}{$\Delta Q+\Delta W=0$}
\end{qoptions}
\qn{6} One mole of an ideal monatomic gas has initial temperature $T_0$, is made to go through the cycle abca shown in the given figure. If $U$ denotes the internal energy, then choose the correct alternative.
\qfig{p1_phy_s1_q06_1}
\begin{qoptions}
\opt{A}{$U_c>U_b>U_a$}
\opt{B}{$U_c-U_b=3RT_0$}
\opt{C}{$U_c-U_a=\dfrac{9RT_0}{2}$}
\opt{D}{$U_b-U_a=\dfrac{3RT_0}{2}$}
\end{qoptions}
\sect{SECTION-I(ii)}
\paraq{7}{8}
On a rough horizontal surface (coefficient of friction $=\mu$) a charged particle of charge $+Q$ and mass $m$ is projected with an initial velocity $v_0$. The charge moves on the surface and eventually comes to rest. Magnetic field $B$ is uniform and perpendicular to the surface as shown.
\qfig{p1_phy_s2_para7_8_1}
\qn{7} Which of the following statement(s) about the velocity of the particle is/are correct :- (as long at it is non zero)
\begin{qoptions}
\opt{A}{The magnitude at any time is $v_0-\mu gt$.}
\opt{B}{The force of magnetic field is perpendicular to it.}
\opt{C}{The velocity vector rotates at a constant angular velocity.}
\opt{D}{The velocity vector rotates at a variable angular velocity.}
\end{qoptions}
\qn{8} Choose the correct statement (s) about the path followed by the particle.
\begin{qoptions}
\opt{A}{The path is circular.}
\opt{B}{The length of the path is $\dfrac{v_0^2}{2\mu g}$.}
\opt{C}{The particle stops at a distance $\dfrac{3v_0^2}{2\mu g}$ from the initial position.}
\opt{D}{The particle stops at a time $\dfrac{v_0}{\mu g}$ after being projected.}
\end{qoptions}
\paraq{9}{10}
Consider the circuit shown in figure. Battery is ideal and capacitor was initially \underline{uncharged}. If at $t=0$ switch is closed, then answer the following question :
\qfig{p1_phy_s2_para9_10_1}
\qn{9} Current in the direction B shown in diagram at $t=0$ (Just after the switch is closed) will be :-
\begin{qoptions}
\opt{A}{$\dfrac38$ A}
\opt{B}{$\left(\dfrac{-3}{8}\right)$ A}
\opt{C}{$\left(\dfrac45\right)$ A}
\opt{D}{$\left(\dfrac{-4}{5}\right)$ A}
\end{qoptions}
\qn{10} Current $I$ as a function of $t$ (in $\mu$s) will be :-
\begin{qoptions}
\opt{A}{$I=\dfrac94e^{-3t/8}$}
\opt{B}{$I=\dfrac32e^{-3t/10}$}
\opt{C}{$I=\dfrac94e^{-3t/16}$}
\opt{D}{$I=\dfrac32e^{-3t/5}$}
\end{qoptions}
\sect{SECTION-II}
\qn{1} A square conducting loop of side length $a=4$m carrying current $i_0=3$A (anticlockwise) is placed in x--y plane. where magnetic field varies as $\vec B=B_0\left(1+\dfrac x4\right)\hat k$ where $B_0=\dfrac18$T. Wire frame is so fixed that it can rotate freely about the central line as shown in figure work done by external agent in joule in rotating this loop slowly by $90^\circ$.
\qfig{p1_phy_s3_q01_1}
\qn{2} A `U' shaped conducting frame is fixed in space. A conducting rod CD lies at rest on the smooth frame. The frame is kept in magnetic field $B$, which is perpendicular to the plane of frame. At time $t=0$, the magnitude of magnetic field begins to change with time $t$ as $B=\dfrac{B_0}{1+5t}$, where $B$ is in tesla, $t$ is in seconds and $B_0$ is positive constant for no current to be even induced in frame the constant speed (in m/s) with which the rod should moved starting from time $t=0$ is
\qfig{p1_phy_s3_q02_1}
\qn{3} An equiconvex lens of refractive index $\dfrac32$ and of radius of curvature 16 cm is placed co-axially at a distance of 90 cm in front of concave mirror of radius of curvature 60 cm. The surrounding medium is air. A point object O is placed on the common axis of the lens and the mirror at a distance 90 cm left of the lens. At an instant ($t=0$) the object, the lens and the mirror start moving with a speed 2 cm/s, 3 cm/s and 8 cm/s respectively in the direction shown in the figure. If the time after which the object and its final image coincide for the second time is $\dfrac{58}{N}$ second then $N$.
\qfig{p1_phy_s3_q03_1}
\qn{4} A parallel plate capacitor with plate of length $\ell$ is included in a circuit as shown in fig. Given are, e.m.f. of the current source, its internal resistance $r$ and the distance $d$ between the plates. As electron with a velocity $u$ flies into the capacitor, so that the electron flies out of the capacitor at an angle $37^\circ$ to the plates. Find the value of $R$ in ohm (assuming circuit in steady state and neglect relativistic effect). Take $\ell=91$ cm, $E=3$V, $r=2\Omega$, $d=\dfrac13$ mm, $u=4\times10^7$ m/s, $M_e=9.1\times10^{-31}$ kg and $e=1.6\times10^{-19}$C.
\qfig{p1_phy_s3_q04_1}
\qn{5} A small conducting square loop of side length $d=10$ cm, is at height $h$ above the boundary of a region of uniform magnetic field $B_0=1$T (into the plane of the paper as shown in figure). The mass of the square loop is $m=100$ gram and resistance $R=0.1\ \Omega$. The loop is dropped from such position and after some time it enters the magnetic field region. Find the height `h' in metre so that velocity of the loop will remain constant during the time interval the loop will be entering into the magnetic field region ($g=10$ m/s$^2$)
\qfig{p1_phy_s3_q05_1}
\qn{6} Two metals bars are fixed vertically and are connected on the top by a capacitor $C$. A sliding conductor AB of length $L$ slides without friction with its ends in contact with the bars. The arrangement is placed in a uniform horizontal magnetic field directed normal to the plane of the figure. The conductor is released from rest. Find the displacement ($x$) in meter of the conductor at time $t=2$ sec.\par
[Given, $M=100$ gm, $g=10$ m/s$^2$, $B=100$ Tesla, $L=1$m and $C=10\ \mu$F]
\qfig{p1_phy_s3_q06_1}
\qn{7} In the figure three process on a gas are shown. The value of $C_V$ is $\alpha R$ for the gas. If AC represents adiabatic process and BA represents isothermal process then $\alpha$ ?
\qfig{p1_phy_s3_q07_1}
\qn{8} A resistanceless rod of length 3m is being pulled in a uniform vertical magnetic field of magnitude 0.8 T with a speed of 5 m/s. The resistance of circuit is 4$\Omega$ and capacitance is 6$\mu$F. Initially there is no charge on the capacitor. What is the current (in A) in circuit when charge on the capacitor becomes 24$\mu$C ?
\qfig{p1_phy_s3_q08_1}
\raggedcolumns
\partsec{2}{CHEMISTRY}
\sect{SECTION-I(i)}
\qn{1} Which of the following statement is/are incorrect :
\begin{qoptions}
\opt{A}{Dipositive zinc exhibits paramagnetism due to loss of two electrons from 3d-orbital of neutral atom.}
\opt{B}{The magnetic moment of an atom is related to the number of unpaired electrons in its electronic configuration.}
\opt{C}{Bohr theory can be successfully modified to explain the electronic spectrum of multielectron atoms.}
\opt{D}{The orbital angular momentum of an electron in an atom is given as $\dfrac{nh}{2\pi}$.}
\end{qoptions}
\qn{2} For the reaction : A $\to$ B. Differential rate law is : $\dfrac{-d[A]}{dt}=k[A]^{-1}$. If initial concentration of A is $A_0$M and concentration at time `t' is $A_t$ M then identify correct option (s)
\begin{qoptions}
\opt{A}{$A_t^2=A_0^2-2kt$}
\opt{B}{Half life, $t_{1/2}=\dfrac{3A_0^2}{4K}$}
\opt{C}{Completion time for reaction : $t_{100\%}=\dfrac{A_0^2}{2K}$}
\opt{D}{The reaction must be a complex reaction}
\end{qoptions}
\qn{3} Which of the following reaction will generate 3$^\circ$ alcohol?
\begin{qoptions}
\opt{A}{\qfigi{p1_chem_s1_q03_A}}
\opt{B}{\qfigi{p1_chem_s1_q03_B}}
\opt{C}{\qfigi{p1_chem_s1_q03_C}}
\opt{D}{\qfigi{p1_chem_s1_q03_D}}
\end{qoptions}
\qn{4} Select reaction(s) correctly matched with the major product :
\begin{qoptions}
\opt{A}{OHC--CHO $\xrightarrow[\Delta]{\text{Conc. NaOH}}$ HOCH$_2$--COO$^-$Na$^+$}
\opt{B}{\qfig{p1_chem_s1_q04_B}\qfig{p1_chem_s1_q04_3}}
\opt{C}{Ph--CO--CO--Ph $\xrightarrow[\text{(ii) H}^+]{\text{(i) OH}^-}$ Ph--C(OH)(Ph)--COOH}
\opt{D}{CH$_3$--CH=O $+$ Br--CH$_2$COOEt $\xrightarrow[\text{(ii) H}_3\text{O}^+\text{ (excess)}]{\text{(i) Zn / ether}}$ CH$_3$--CH(OH)--CH$_2$--COOH}
\end{qoptions}
\qn{5} Which of the following statement(s) regarding the pair of complexes is/are incorrect?
\begin{qoptions}
\opt{A}{Both K$_2$[PdCl$_4$] and K$_2$[Pt(ox)$_3$] are optically inactive}
\opt{B}{Both K$_2$[PdCl$_4$] and K$_2$[Pt(ox)$_3$] do not from diastereomers.}
\opt{C}{K$_2$[Ni(CN)$_4$] and Co[Hg(SCN)$_4$] have same hybridization of central atom.}
\opt{D}{K$_2$[Ni(CN)$_4$] and Co[Hg(SCN)$_4$] both can show same type of structural isomerism and are diamagnetic.}
\end{qoptions}
\qn{6} Select similarity between Br$^-$ and I$^-$ ion :-
\begin{qoptions}
\opt{A}{Both can be oxidised by conc. H$_2$SO$_4$ to produce respective halogen}
\opt{B}{Both can be oxidised by Cl$_2$ water to produce respective halogen}
\opt{C}{Both can be oxidised by Cu$^{2+}_{\text{(aq.)}}$ to produce respective halogen}
\opt{D}{Both are precipitated by AgNO$_3$ solution}
\end{qoptions}
\sect{SECTION-I(ii)}
\paraq{7}{8}
\noindent(A)$_{\text{aq}}$ $\xrightarrow{\text{(Zn + dil. H}_2\text{SO}_4)}$ (B) $\xrightarrow{\text{FeCl}_3\,/\,\text{H}^+}$ (D) $+$ (E) (Coloured ion)\par
(A)$_{\text{aq}}$ $\xrightarrow{\text{(D)}}$ (F)\par
(B) $\xrightarrow{\text{Cr}_2\text{O}_7^{2-}/\text{H}^+}$ (D) (White turbidity) $+$ Cr$^{+3}_{\text{(aq.)}}$
\qn{7} Species (A) and (F) are respectively :
\begin{qoptions}
\opt{A}{SO$_{3(\text{aq.})}^{-2}$, S}
\opt{B}{SO$_{3(\text{aq.})}^{-2}$, S$_2$O$_{3(\text{aq.})}^{-2}$}
\opt{C}{S$_2$O$_{3(\text{aq.})}^{-2}$, SO$_{3(\text{aq.})}^{-2}$}
\opt{D}{SO$_{3(\text{aq.})}^{2-}$, $\bar{\text{S}}$CN$_{\text{aq.}}$}
\end{qoptions}
\qn{8} When FeCl$_3$ react with NH$_4$SCN then :
\begin{qoptions}
\opt{A}{Blood red complex is formed}
\opt{B}{Blue coloured complex is formed}
\opt{C}{Brown coloured complex is formed}
\opt{D}{Green solution of Fe$^{2+}$}
\end{qoptions}
\paraq{9}{10}
Consider a three step mechanism\par
\noindent(i) I$_2$ $\underset{k_{-1}}{\overset{k_1}{\rightleftharpoons}}$ 2I \quad (fast)\par
\noindent(ii) H$_2$ $+$ I $\underset{k_{-2}}{\overset{k_2}{\rightleftharpoons}}$ H$_2$I \quad (fast)\par
\noindent(iii) H$_2$I $+$ I $\xrightarrow{k_3}$ 2HI \quad (slowest)
\qn{9} Which is/are incorrect regarding above mechanism ?
\begin{qoptions}
\opt{A}{Rate of reaction $=K[\text{H}_2\text{I}]\,[\text{I}]$ where $K=k_3\cdot\dfrac{k_2}{k_{-2}}\cdot\left(\dfrac{k_1}{k_{-1}}\right)$}
\opt{B}{Rate is decided only by the steps (i) and (ii)}
\opt{C}{Molecularity of overall reaction is 2}
\opt{D}{Order of overall reaction is 2.}
\end{qoptions}
\qn{10} Which of the following qualitative plot of potential energy (PE) v/s reaction coordinates is/are incorrect if $k_1>>k_2>>k_3$ and reaction is exothermic :
\begin{qoptions}
\opt{A}{\qfigi{p1_chem_s2_q10_A}}
\opt{B}{\qfigi{p1_chem_s2_q10_B}}
\opt{C}{\qfigi{p1_chem_s2_q10_C}}
\opt{D}{\qfigi{p1_chem_s2_q10_D}}
\end{qoptions}
\sect{SECTION-II}
\qn{1} A sample of hydrogen atoms contains few atoms in two different excited state $n_1$ and $n_2$ only. These atoms make transition to ground state. If the energy of spectral line obtained due to transition of electron from $n_2$ to $n_1$ is 2.55 eV then calculate the maximum number of different lines obtained in spectrum.
\qn{2} How many of the following statement(s) are true (Assuming all gases involved to behave ideally) ?\par
(i) H$_2$O($\ell$) $\longrightarrow$ H$_2$O(g) \quad (1 atm, 373 K $\to$ 1 atm, 373 K)\par
For this $\Delta H=P\Delta V$ as $\Delta U=0$\par
(ii) For the ideal gas reaction H$_2$(g) $+$ I$_2$(g) $\longrightarrow$ 2HI(g) occurring in a closed rigid container $\Delta H=q_v$ where $q_v$ is the heat absorbed/released at constant volume.\par
(iii) For any irreversible isothermal expansion of an ideal gas $|\Delta S_{\text{system}}|>|\Delta S_{\text{surr}}|$\par
(iv) For an ideal gas $\Delta H=nC_{p,m}\Delta T$ is applicable only under constant pressure condition for temperature change $\Delta T$.\par
(v) The change in enthalpy of the gas is (a) zero if it is expanded reversibly with $T_i=T_f$ and (b) positive if it is expanded reversibly under adiabatic condition with $T_i\neq T_f$.\par
(vi) If the expansion is carried out freely, no change in internal energy, enthalpy and entropy is observed.\par
(vii) P (black, solid) $\longrightarrow$ P (white, solid) $\Delta S=$ positive\par
(viii) For mixing equal volumes of two ideal gases at constant temperature and pressure in an isolated container, $\Delta S_{\text{system}}>0$ and $\Delta S_{\text{total}}>0$
\qn{3} Total number of hydroxyl groups present in a molecule of the major product P is \underline{\hspace{1cm}}
\qfig{p1_chem_s3_q03_1}
\qn{4} The total number of all possible Mono Aldols ($\beta$-Hydroxy carbonyl compounds) are formed in the following reaction. (Ignore stereoisomers)\par
PhCHO $+$ CH$_3$CH=O $+$ Ph--CO--CH$_3$ $\xrightarrow{\text{dil. NaOH, }\Delta}$
\qfullbegin
\qn{5} How many compounds gives iodoform test with NaOH $+$ I$_2$ :
\qfig{p1_chem_s3_q05_1}
\qfullend
\qn{6} Select the complexes which will absorb in visible range :\par
(a) [Ti(H$_2$O)$_6$]Cl$_3$\par
(b) [Cu(NH$_3$)$_4$]$^+$\par
(c) [Cu(NH$_3$)$_4$]$^{2+}$\par
(d) [Ni(en)$_3$]SO$_4$\par
(e) [Ni(CO)$_4$]\par
Hence, write total number of such complexes :
\qn{7} Identify total number of compounds which liberate `CH$_4$' gas on reaction with CH$_3$MgCl.\par
(i) H$_2$O \quad (ii) ethene \quad (iii) Me--C$\equiv$C--H \quad (iv) ethane\par
(v) Ph--OH \quad (vi) CCl$_4$ \quad (vii) MeCOOH
\qn{8} Find sum of the oxidation number of metal in products of given reactions.\par
(i) Cu$^{2+}$ $+$ S$_2$O$_3^{2-}$ $_{\text{(excess)}}$ $\to$\par
(ii) Fool's gold $+$ O$_2$ $_{\text{(excess)}}$ $\xrightarrow{\text{roasting}}$\par
(iii) KCl $+$ MnO$_2$ $\xrightarrow{\text{H}^+/\Delta}$
\raggedcolumns
\partsec{3}{MATHEMATICS}
\sect{SECTION-I(i)}
\qn{1} Let $f(x)=\cot^{-1}\left(x^2+4x+\alpha^2-3\alpha\right)$ defined on $R\to\left(0,\dfrac\pi2\right]$, is an onto function then
\begin{qoptions}
\opt{A}{$\alpha\in[-1,4]$}
\opt{B}{$f'(0)=-4/17$}
\opt{C}{$f(x)$ is one-one function}
\opt{D}{$f(x)$ is many-one function}
\end{qoptions}
\qn{2} Let $f:R\to R$ be defined by $f(x)=\begin{cases}x+2x^2\sin\dfrac1x & \text{for } x\neq0\\ 0 & \text{for } x=0\end{cases}$ then
\begin{qoptions}
\opt{A}{$f(x)$ is differentiable for all $x$, but $f'(x)$ is not continuous at $x=0$}
\opt{B}{$f'(0)=1$}
\opt{C}{$f(x)$ is increasing at $x=0$}
\opt{D}{Both $f(x)$ and $f'(x)$ are differentiable for all $x$.}
\end{qoptions}
\qn{3} Area bounded by curves $y=\sqrt{|x-1|}$ and $\sqrt y=|x-1|$ is
\begin{qoptions}
\opt{A}{$2\displaystyle\int_1^2\left(\sqrt{x-1}-(x-1)^2\right)dx$}
\opt{B}{$2\displaystyle\int_0^1\left(\sqrt x-x^2\right)dx$}
\opt{C}{$2\displaystyle\int_1^2\left(\sqrt{2-x}-(2-x)^2\right)dx$}
\opt{D}{$\dfrac23$}
\end{qoptions}
\qn{4} Let $f(x)=x\ln\left(1+\dfrac1x\right)$, $0<x<\infty$, then which of the following is correct
\begin{qoptions}
\opt{A}{$f$ is strictly increasing}
\opt{B}{$\displaystyle\lim_{x\to0}f(x)=0$}
\opt{C}{$\displaystyle\lim_{x\to\infty}f(x)=1$}
\opt{D}{$f$ has exactly one point of extremum}
\end{qoptions}
\qn{5} Let $f(x)=2x^3-15x^2+36x-23$ and $g(x)=\begin{cases}\max\{f(t),\,1\le t\le x\}, & \text{if } 1\le x\le\dfrac72\\ 12-2x, & \text{if } \dfrac72<x\le6\end{cases}$. Which of the following is correct about $g(x)$?
\begin{qoptions}
\opt{A}{Continuous everywhere}
\opt{B}{Differentiable everywhere}
\opt{C}{has greatest value 5}
\opt{D}{has least value 0}
\end{qoptions}
\qn{6} Let $C$ be a curve such that normal at any point $P$ on it meets x-axis and y-axis at $A$ and $B$ respectively. If $BP:PA=1:2$ (internally) and curve passes through $(0,4)$, then-
\begin{qoptions}
\opt{A}{The curve `C' passes through $(\sqrt{10},-6)$}
\opt{B}{Equation of tangent at $(4,4\sqrt3)$ is $\sqrt3x-2y+4\sqrt3=0$}
\opt{C}{Equation of tangent at $(4,4\sqrt3)$ is $2x-\sqrt3y+4=0$}
\opt{D}{If point $P$ is $(8,12)$, then point $A$ and point $B$ are $(24,0)$ and $(0,18)$ respectively.}
\end{qoptions}
\sect{SECTION-I(ii)}
\paraq{7}{8}
For any geometrical figure possessing a finite area `A' and a finite perimeter `P', the ratio $\dfrac{4\pi A}{P^2}$ is called the isoperimetric quotient (I.Q.)
\qn{7} If parallelogram ABCD and the triangle ABC have the same I.Q., then the ratio $\dfrac{AB+BC}{AC}$ is equal to :
\begin{qoptions}
\opt{A}{$1+2\sqrt2$}
\opt{B}{$\dfrac{1}{\sqrt2-1}$}
\opt{C}{$\dfrac{4+\sqrt2}{2}$}
\opt{D}{$1+\sqrt2$}
\end{qoptions}
\qn{8} A sector of a circle has a fixed radius $r$ and a variable angle $\theta$ (in radians). The value of $\theta$ for which I.Q. of the sector is maximum is :
\begin{qoptions}
\opt{A}{2}
\opt{B}{$\dfrac\pi2$}
\opt{C}{1}
\opt{D}{An integer}
\end{qoptions}
\paraq{9}{10}
Suppose that function $f:R\to R$ satisfies the inequality, $\left|\displaystyle\sum_{k=1}^{n}\left(3^k\{f(x+ky)-f(x-ky)\}\right)\right|\le1$ for every positive integer $n$ and for all $x,y\in R$. Also, let $g(x)=ax^3-(3+b)x^2+cx$ for $x\in R$, where $a$, $b$, $c$ are such that $\displaystyle\lim_{x\to0}\dfrac{axe^x-b\ln(1+x)+cxe^{-x}}{x^2\sin x}=2$
\qn{9} If $f(13)=-23$, number of solutions of $g(x)=f(x)$ is more than or equal to :
\begin{qoptions}
\opt{A}{2}
\opt{B}{3}
\opt{C}{4}
\opt{D}{1}
\end{qoptions}
\qn{10} The value of $\displaystyle\lim_{x\to1}\dfrac{g(x)+3}{\ln(2-x)}$ is equal to
\begin{qoptions}
\opt{A}{12}
\opt{B}{3}
\opt{C}{9}
\opt{D}{1}
\end{qoptions}
\sect{SECTION-II}
\qn{1} Let parabola $P$ passes through the points $(0,1)$, $(2,0)$, $(0,2)$ and $(2,2)$. If axis of symmetry of the parabola has slope $m$, then $|m|$ is equal to \underline{\hspace{1.5cm}}.
\qn{2} A ray emanating from the point $(0,6)$ is incident on the ellipse $25x^2+16y^2=1600$ at the point $P$ with ordinate 5. After reflection, ray cuts the Y-axis at $B$. The length of $PB/2$ is equal to \underline{\hspace{1.5cm}}.
\qn{3} The value of the integral $\displaystyle\int_0^1\dfrac{\sin^{-1}\sqrt x}{x^2-x+1}\,dx$ is equal to \underline{\hspace{1.5cm}}.
\qn{4} Let $f(x)$ be a polynomial with leading coefficient unity and satisfying the equation $(x-16)f(2x)=16(x-1)f(x)$ then $\dfrac{f(0)}{128}$ is equal to \underline{\hspace{1.5cm}}.
\qn{5} $f:[-1,1]\to\left[-\dfrac12,\dfrac12\right]$, $f(x)=\dfrac{x}{1+x^2}$. If the area bounded by $y=f^{-1}(x)$, X-axis between, $x=\dfrac{-1}{2}$ to $x=\dfrac12$ is $A-\ln2$, then the value of $A$ is equal to \underline{\hspace{1.5cm}}.
\qn{6} $f(x)$ is a fifth order polynomial in $x$ with every root of $f(x)=0$ is real and distinct, then the number of real roots of $f''(x)\,f(x)-(f'(x))^2=0$ is equal to \underline{\hspace{1.5cm}}.
\qn{7} If $\displaystyle\int(x^9+x^6+x^3)(2x^6+3x^3+6)^{1/3}dx=a\left(2x^9+3x^6+6x^3\right)^{4/3}+C$, then the value of $48a$ is equal to \underline{\hspace{1.5cm}}.\par
(Where $C$ is constant of integration)
\qn{8} If $\displaystyle\int\dfrac{\sec x\,(2+\sec x)}{(1+2\sec x)^2}\,dx$ is $\left(\dfrac{\cos x+2}{\lambda\sin x}\right)^p$, then $\lambda+p$ is equal to \underline{\hspace{1.5cm}}.
\raggedcolumns
\newpage
\papersec{2}
\partsec{1}{PHYSICS}
\sect{SECTION-I}
\qn{1} A circuit diagram is shown in figure.
\qfig{p2_phy_s1_q01_1}
Which of the following option is incorrect.
\begin{qoptions}
\opt{A}{Current flowing in resistance R is zero.}
\opt{B}{Terminal voltage of 10 V battery is zero.}
\opt{C}{Terminal voltage of 5 V battery is zero.}
\opt{D}{Current flowing through battery of 10 V is zero.}
\end{qoptions}
\qn{2} A charged particle of mass 1 kg and charge 1 C moves with velocity $\vec v=6\hat i+8\hat j$ m/s in a magnetic field of $\vec B=2\hat k$ T. Then :
\begin{qoptions}
\opt{A}{the path of particle may be $y^2+z^2=25$}
\opt{B}{the path of particle may be $x^2+z^2=25$}
\opt{C}{the path of particle may be $x^2+y^2=25$}
\opt{D}{time period of particle will be 2s.}
\end{qoptions}
\qn{3} Three resistances of magnitude $R$ each are connected in the form of an equilateral triangle of side $a$. The combination is placed in a magnetic field $B=B_0\,e^{-\lambda t}$ perpendicular to its plane. The induced current in the circuit is given by
\qfig{p2_phy_s1_q03_1}
\begin{qoptions}
\opt{A}{$\left(\dfrac{a^2\lambda}{2\sqrt3R}B_0\right)e^{-\lambda t}$}
\opt{B}{$\left(\dfrac{a^2\lambda}{4\sqrt3R}B_0\right)e^{-\lambda t}$}
\opt{C}{$\left(\dfrac{a^2B_0}{\lambda4\sqrt3R}\right)e^{-\lambda t}$}
\opt{D}{$\left(\dfrac{a^2B_0R}{\lambda4\sqrt3}\right)e^{-\lambda t}$}
\end{qoptions}
\qn{4} Two convex lenses of focal length 10 cm and 20 cm respectively placed coaxially and are separated by some distance `d'. The whole system behaves like a concave lens for rays coming from infinity. One of the possible value of `d' is :
\begin{qoptions}
\opt{A}{40 cm}
\opt{B}{25 cm}
\opt{C}{15 cm}
\opt{D}{20 cm}
\end{qoptions}
\qn{5} Four charges of 8$\mu$C, $-6\mu$C, 2 $\mu$C and $-4\mu$C are placed at corners of a square placed in xy plane with its center at origin. The electric potential due do these charges is zero every where on the line
\begin{qoptions}
\opt{A}{$x=y$, $z=0$}
\opt{B}{$x=0=z$}
\opt{C}{$x=0=y$}
\opt{D}{$x=z$, $y=0$}
\end{qoptions}
\qn{6} A particle starts oscillating simple harmonically from its extreme position at $t=0$. The time period of oscillation is $T$. The ratio of kinetic energy and potential energy of particle at $t=\dfrac T6$ is
\begin{qoptions}
\opt{A}{2 : 1}
\opt{B}{3 : 1}
\opt{C}{4 : 1}
\opt{D}{1 : 4}
\end{qoptions}
\qn{7} Figure shows a parabolic curve between $\dfrac1V$ and $T$ for a mixture of a gas undergoing adiabatic process.
\qfig{p2_phy_s1_q07_1}
The value of adiabatic exponent ($\gamma$) of mixture is
\begin{qoptions}
\opt{A}{$\dfrac32$}
\opt{B}{$\dfrac53$}
\opt{C}{$\dfrac75$}
\opt{D}{$\dfrac52$}
\end{qoptions}
\qn{8} In the circuit shown in figure.
\qfig{p2_phy_s1_q08_1}
\begin{qoptions}
\opt{A}{$q_2=280\mu$C}
\opt{B}{$q_2=360\mu$C}
\opt{C}{$q_3=360\mu$C}
\opt{D}{It is not possible to find $q_2$ and $q_3$ unless $C_1$ and $C_2$ are known.}
\end{qoptions}
\qn{9} Block of mass $m=4$ kg as shown in figure is slightly displaced and released. The spring constant $k=100$ N/m. Pulleys are massless and frictionless. The time period of oscillation of block will be :
\qfig{p2_phy_s1_q09_1}
\begin{qoptions}
\opt{A}{$\dfrac\pi2$ s}
\opt{B}{$\dfrac\pi4$ s}
\opt{C}{$\dfrac{\pi}{10}$ s}
\opt{D}{$\dfrac{2\pi}{3}$ s}
\end{qoptions}
\qn{10} Two different perfect monoatomic gases at absolute temperature 300 K and 200 K having number of molecules $2\times10^{18}$ and $3\times10^{18}$ are mixed in a container of negligible heat capacity. Temperature of mixture is :
\begin{qoptions}
\opt{A}{210 K}
\opt{B}{240 K}
\opt{C}{260 K}
\opt{D}{280 K}
\end{qoptions}
\qn{11} 100 J of heat is given to a monatomic gas sample slowly such that pressure remains constant. The work done by gas is :
\begin{qoptions}
\opt{A}{40 J}
\opt{B}{80 J}
\opt{C}{60 J}
\opt{D}{20 J}
\end{qoptions}
\qn{12} A resistance is shown in the figure. Its value (in $\Omega$) and tolerance are given respectively.
\qfig{p2_phy_s1_q12_1}
\begin{qoptions}
\opt{A}{$14\times10^5$, 5 \%}
\opt{B}{$14\times10^5$, 10 \%}
\opt{C}{$27\times10^5$, 10\%}
\opt{D}{$27\times10^5$, 20\%}
\end{qoptions}
\qn{13} Four simple harmonic motions\par
$x_1=8\sin100\,\pi t$\par
$x_2=6\sin\left(100\,\pi t+\dfrac\pi2\right)$\par
$x_3=4\sin(100\,\pi t+\pi)$\par
$x_4=2\sin\left(100\,\pi t+\dfrac{3\pi}2\right)$\par
are superimposed on each other. The resulting amplitude of vibration is :
\begin{qoptions}
\opt{A}{$4\sqrt2$ unit}
\opt{B}{10 unit}
\opt{C}{4 unit}
\opt{D}{20 unit}
\end{qoptions}
\qn{14} A magnetic needle is free to rotate in vertical plane which makes an angle of $60^\circ$ with magnetic meridian. If the needle stays in a direction making an angle of $\tan^{-1}\left(\dfrac{2}{\sqrt3}\right)$ with horizontal. What would be the dip at that place.
\begin{qoptions}
\opt{A}{$37^\circ$}
\opt{B}{$30^\circ$}
\opt{C}{$45^\circ$}
\opt{D}{$53^\circ$}
\end{qoptions}
\qn{15} Voltage versus time and current versus time graph of a circuit element is shown in figure. The type of circuit element is :
\qfig{p2_phy_s1_q15_1}
\begin{qoptions}
\opt{A}{resistance of 2$\Omega$}
\opt{B}{capacitance of 1F}
\opt{C}{capacitance of 2F}
\opt{D}{none}
\end{qoptions}
\qn{16} Four long straight long wire are leaked at corners of square PQRS. All the wires carry equal currents. Currents in wire P and Q are outwards and in R and S are inwards. The magnetic field at O is along :
\qfig{p2_phy_s1_q16_1}
\begin{qoptions}
\opt{A}{PQ}
\opt{B}{QR}
\opt{C}{SP}
\opt{D}{RS}
\end{qoptions}
\qn{17} One mole of an ideal monatomic gas undergoes the process $P=2\sqrt T$. The molar specific heat of gas is
\begin{qoptions}
\opt{A}{$\dfrac{3R}{4}$}
\opt{B}{$\dfrac{9R}{4}$}
\opt{C}{2R}
\opt{D}{$\dfrac{2R}{3}$}
\end{qoptions}
\qn{18} A charged particle enters in a uniform magnetic field with velocity vector at an angle $30^\circ$ with field. The pitch of helical path is $p$. The radius of helix will be :
\begin{qoptions}
\opt{A}{$\dfrac{p}{2\pi}$}
\opt{B}{$\dfrac{p}{2\sqrt3\pi}$}
\opt{C}{$\dfrac{2\pi p}{3}$}
\opt{D}{$\dfrac{p}{\sqrt3\pi}$}
\end{qoptions}
\qn{19} A heat pump operating between reservoirs at $-5^\circ$C and $20^\circ$C requires an electrical input of 1.2 kJ in each cycle. The heat dissipated into hot reservoir in each cycle.
\begin{qoptions}
\opt{A}{12.8 kJ}
\opt{B}{11.7 kJ}
\opt{C}{14 kJ}
\opt{D}{10 kJ}
\end{qoptions}
\qn{20} A charged particle of mass `m' and charge $q$ is moving directly towards a fixed charge $Q$. When it is at a distance `d' from fixed charge it has velocity $v$. The distance of closest approach will be. (Here, $\dfrac{\pi\epsilon_0mv^2}{qQ}=\dfrac2d$)
\begin{qoptions}
\opt{A}{$\dfrac d2$}
\opt{B}{$\dfrac d3$}
\opt{C}{$\dfrac d5$}
\opt{D}{$\dfrac{3d}{4}$}
\end{qoptions}
\sect{SECTION-II}
\qn{1} There is a potentiometer wire AB of resistance 100 $\Omega$ and rheostat of variable resistance `$x$' as shown in figure. For $x=0$, null point is found at 30 cm from A. For unknown value of $x$ null point was 40 cm from A. What is the value of $3x$ (in $\Omega$)
\qfig{p2_phy_s2_q01_1}
\qn{2} What amount of charge (in $\mu$C) will flow through switch S from A to B, when switch S is closed.
\qfig{p2_phy_s2_q02_1}
\qn{3} A particle of charge $q$ and mass $m$ starts moving from origin under action of an electric field $\vec E=E_0\hat i$ N/C and magnetic field $\vec B=B_0\hat k$ T. Its velocity at $(x,3,0)$ is $\left(4\hat i+3\hat j\right)$ m/s. The value of $x$ (in m) is $\dfrac{pm}{2qE_0}$. The value of $p$ is
\qn{4} A circular loop of area 10 cm$^2$ and resistance 4$\Omega$ is rotated in a magnetic field $B=2$T through $180^\circ$ in 20s as shown in figure.
\qfig{p2_phy_s2_q04_1}
If magnitude of amount of charge flown in circular loop in that interval is $x\ \mu$C then find value of $x$.
\qn{5} A circuit diagram is shown in figure
\qfig{p2_phy_s2_q05_1}
Find current $I$ (in A) flowing through battery $E$.
\qn{6} An astronomical telescope has angular magnification of magnitude 6 for distant objects. The separation between object and eye piece is 42 cm and final image is formed at infinity. Find focal length of objective (in cm).
\qn{7} The reading of ammeter is 10 A and voltmeter is 100 V as shown in figure.
\qfig{p2_phy_s2_q07_1}
If resistance of voltmeter is 3000 $\Omega$ then find value of resistance `R'.
\qn{8} The electric potential existing in space is given as $V(x,y,z)=10\,(xy+yz+zx)$ in SI unit. Find the value of electric field (in N/C) at point $(2,2,2)$ (approximately).
\qn{9} An ice cube of mass 400 g at temperature $0^\circ$C melts to water at $0^\circ$C. Heat is supplied from the surrounding air at slightly above $0^\circ$C. The process takes place very slowly. Find the change in entropy (J/K) of ice when it has all melted. (Latent heat of ice $=334$ kJ/kg)
\qn{10} A rectangular loop PQRS is formed by joining the corners of cube of side 20 cm as shown in figure. If $I=8$ A and $\vec B=0.4\hat i$ T then the magnitude of torque on loop is $P\times10^{-3}$ N-m. Find the value of $P$.
\qfig{p2_phy_s2_q10_1}
\raggedcolumns
\partsec{2}{CHEMISTRY}
\sect{SECTION-I}
\qn{1} Joseph Priestley discovered sweet smelling, colorless laughing gas, nitrous oxide (N$_2$O), in 1763. N$_2$O has been widely used as an anesthetic for medical and dental applications. Nitrogen dioxide (NO$_2$) is a toxic reddish brown gas that is the major component of smog. Nitrogen oxide (NO) is a colorless, non-toxic gas that readily oxidizes to nitrogen dioxide in the presence of oxygen. Using the reactions :\par
4NO(g) $\to$ 2N$_2$O(g) $+$ O$_2$(g) \quad $\Delta H^0=-139.56$ kJ\par
2NO(g) $+$ O$_2$(g) $\to$ 2NO$_2$(g) \quad $\Delta H^0=-69.70$ kJ\par
Calculate the standard enthalpy change when 2 moles of nitrous oxide combine with three moles of oxygen to produce four moles of nitrogen dioxide.
\begin{qoptions}
\opt{A}{$-209.26$ kJ}
\opt{B}{$-139.86$ kJ}
\opt{C}{$-69.86$ kJ}
\opt{D}{$+0.16$ kJ}
\end{qoptions}
\qn{2} Which of the following statements is correct :
\begin{qoptions}
\opt{A}{All covalent solids are good conductors of electricity}
\opt{B}{In rhombohedral crystal system, only primitive unit cell is possible}
\opt{C}{Ionic solids are ductile and brittle.}
\opt{D}{In molecular solids, constituted particles are held together through covalent bond.}
\end{qoptions}
\qn{3} For an exothermic reaction whose mechanism is given as :\par
A $+$ B $\to$ X \ (slow)\par
X $\to$ AB \ (fast)\par
The progress of the reaction can be best described by :
\begin{qoptions}
\opt{A}{\qfigi{p2_chem_s1_q03_A}}
\opt{B}{\qfigi{p2_chem_s1_q03_B}}
\opt{C}{\qfigi{p2_chem_s1_q03_C}}
\opt{D}{\qfigi{p2_chem_s1_q03_D}}
\end{qoptions}
\qn{4} From the following observations, predict the type of orbital.\par
Observation 1 : xy-plane acts as nodal plane.\par
Observation 2 : The angular function of the orbital intersect the three axis at origin only.\par
Observation 3 : R$^2$(r) v/s r curve obtained for the orbital is :
\qfig{p2_chem_s1_q04_1}
\begin{qoptions}
\opt{A}{5p$_z$}
\opt{B}{6d$_{xy}$}
\opt{C}{6d$_{x^2-y^2}$}
\opt{D}{6d$_{yz}$}
\end{qoptions}
\qn{5} When 5 mL of a 1.0 M HCl solution is mixed with 5 mL of a 0.1 M NaOH solution, temperature of solution increases by $t^\circ$C. Which of the following can be predicted from this observation :
\begin{qoptions}
\opt{A}{If 5 mL of 0.1M HCl is mixed with 5 mL 0.1 M NH$_3$ solution, the temperature rise will be more than $t^\circ$C}
\opt{B}{If 5 mL of 0.1 M CH$_3$COOH is mixed with 5 mL 0.1 M NaOH, the temperature rise will be less than $t^\circ$C}
\opt{C}{If 10 mL of 1.0 M HCl mixed with 10 mL of 0.1 M NaOH, the temperature rise will be $2t^\circ$C}
\opt{D}{If 10 mL of 0.1 M HCl is mixed with 10 mL of 0.1 M NaCl, the temperature rise will be $t^\circ$C}
\end{qoptions}
\qn{6} Predict the major organic product of the following reactions.\par
(CH$_3$)$_2$CHCO$_2$H $\xrightarrow[\text{(ii) H}_2\text{O}]{\text{(i) LiAlH}_4}$ $\xrightarrow{\text{PBr}_3}$ $\xrightarrow[\text{DMSO}]{\text{KCN}}$ $\xrightarrow{\text{H}_2\text{O, H}^+}$
\begin{qoptions}
\opt{A}{(CH$_3$)$_2$CHCH$_2$CH$_2$NH$_2$}
\opt{B}{(CH$_3$)$_2$CHCHBrCO$_2$H}
\opt{C}{(CH$_3$)$_2$C $=$ CHCO$_2$H}
\opt{D}{(CH$_3$)$_2$CHCH$_2$CO$_2$H}
\end{qoptions}
\qn{7} Which of the following reaction produces the same product :
\qfig{p2_chem_s1_q07_1}
\qfig{p2_chem_s1_q07_2}
\begin{qoptions}
\opt{A}{I, \& II}
\opt{B}{I, II, III \& IV}
\opt{C}{I, III \& IV}
\opt{D}{II, III \& IV}
\end{qoptions}
\qn{8} Which of the following is incorrectly matched with its major product :
\begin{qoptions}
\opt{A}{\qfigi{p2_chem_s1_q08_A}}
\opt{B}{\qfigi{p2_chem_s1_q08_B}}
\opt{C}{\qfigi{p2_chem_s1_q08_C}}
\opt{D}{\qfigi{p2_chem_s1_q08_D}}
\end{qoptions}
\qn{9} End product of the following sequence of reactions is :
\qfig{p2_chem_s1_q09_1}
\begin{qoptions}
\opt{A}{\qfigi{p2_chem_s1_q09_A}}
\opt{B}{\qfigi{p2_chem_s1_q09_B}}
\opt{C}{\qfigi{p2_chem_s1_q09_C}}
\opt{D}{\qfigi{p2_chem_s1_q09_D}}
\end{qoptions}
\qn{10} The final products of the given reaction is :
\qfig{p2_chem_s1_q10_1}
\begin{qoptions}
\opt{A}{\qfigi{p2_chem_s1_q10_A}}
\opt{B}{\qfigi{p2_chem_s1_q10_B}}
\opt{C}{\qfigi{p2_chem_s1_q10_C}}
\opt{D}{\qfigi{p2_chem_s1_q10_D}}
\end{qoptions}
\qn{11} The major product obtained in the following reaction is :
\qfig{p2_chem_s1_q11_1}
\begin{qoptions}
\opt{A}{\qfigi{p2_chem_s1_q11_A}}
\opt{B}{\qfigi{p2_chem_s1_q11_B}}
\opt{C}{\qfigi{p2_chem_s1_q11_C}}
\opt{D}{\qfigi{p2_chem_s1_q11_D}}
\end{qoptions}
\qn{12}
\qfig{p2_chem_s1_q12_1}
Final product in above reaction sequence?
\begin{qoptions}
\opt{A}{\qfigi{p2_chem_s1_q12_A}}
\opt{B}{\qfigi{p2_chem_s1_q12_B}}
\opt{C}{\qfigi{p2_chem_s1_q12_C}}
\opt{D}{\qfigi{p2_chem_s1_q12_D}}
\end{qoptions}
\qn{13} XeF$_n$ $+$ mH$_2$O $\to$ A (Contain Xe) $+$ Other Products\par
Regarding above reaction, which of the following option is correct :
\begin{qoptions}
\opt{A}{When $n=6$ and $m=2$, the hybridisation of Xe in A is sp$^3$d$^2$.}
\opt{B}{When $n=2$ and $m=1$, then other products has a diatomic colored gas.}
\opt{C}{When $n=6$ and $m=3$, then Electron pair structure of A is tetrahedral.}
\opt{D}{When $n=6$ and $m=1$, the hybridisation of Xe in A is sp$^3$d.}
\end{qoptions}
\qn{14} Select the \textbf{CORRECT} option, for mentioned properties
\begin{qoptions}
\opt{A}{Ba$^{2+}$ $>$ Ca$^{2+}$ $>$ Mg$^{2+}$ $>$ Be$^{2+}$ : (polarising power of cation)}
\opt{B}{BeSO$_4$ $>$ MgSO$_4$ $>$ CaSO$_4$ $>$ SrSO$_4$ : (thermal stability)}
\opt{C}{BF$_3$ $<$ BCl$_3$ $<$ BBr$_3$ $<$ BI$_3$ : (lewis acidic nature)}
\opt{D}{O$_2$ $<$ KO$_2$ $<$ K$_2$O$_2$ : (O--O bond order)}
\end{qoptions}
\qn{15} For [Fe(C$_2$O$_4$)$_3$]$^{3-}$, EAN of Fe will be :
\begin{qoptions}
\opt{A}{27}
\opt{B}{24}
\opt{C}{35}
\opt{D}{29}
\end{qoptions}
\qn{16} Which of the following compounds produce NH$_3$ on heating :
\begin{qoptions}
\opt{A}{NH$_4$Cl}
\opt{B}{(NH$_4$)$_2$S}
\opt{C}{(NH$_4$)$_2$CO$_3$}
\opt{D}{All of these}
\end{qoptions}
\qn{17}
\qfig{p2_chem_s1_q17_1}
\begin{qoptions}
\opt{A}{O$_3$}
\opt{B}{H$_2$S}
\opt{C}{SO$_2$}
\opt{D}{NH$_3$}
\end{qoptions}
\qn{18} For a first order gaseous reaction half life is 25 minutes and reaction started with 1 atm pressure of pure A, then find final pressure of reaction mixture after 75 minutes.\par
A(g) $\to$ 4B(g) $+$ C(g)
\begin{qoptions}
\opt{A}{2 atm}
\opt{B}{3 atm}
\opt{C}{3.75 atm}
\opt{D}{4.5 atm}
\end{qoptions}
\qn{19} Two mole of an ideal gas is isothermally compressed in two stages starting from (1 atm, 20 lit) to 10 lit. by constant external pressure $P_1$ then upto 2 lit by constant external pressure $P_2$. Find work done in the process.
\begin{qoptions}
\opt{A}{80 lit-atm}
\opt{B}{100 lit-atm}
\opt{C}{100J}
\opt{D}{180 lit-atm}
\end{qoptions}
\qn{20} [Ag(CN)$_2$]$^-$ is formed :-
\begin{qoptions}
\opt{A}{In leaching of silver from Ag$_2$S using NaCN$_{\text{(aq.)}}$ \& O$_2$}
\opt{B}{In leaching of silver from native Ag using NaCN$_{\text{(aq.)}}$ \& O$_2$}
\opt{C}{When ppt. of AgCl reacts with excess NaCN}
\opt{D}{All of the above}
\end{qoptions}
\sect{SECTION-II}
\qn{1} Consider the following wave representing electromagnetic radiation :
\qfig{p2_chem_s2_q01_1}
Calculate the energy of above wave (in eV) :
\qn{2} An ideal solution was prepared by dissolving some amount of cane sugar (non-volatile) in 0.9 moles of water. The solution was then cooled just below its freezing temperature ($T_f=271$ K), where some ice get separated out. The remaining aqueous solution registered a vapour pressure of 700 torr at 373 K. Calculate the mass (in gm) of ice separated out ($K_f$ of water $=1.86$ K.kg.mol$^{-1}$)
\qn{3} Total number of monobrominated product obtained in the given reaction is :\par
CH$_2$=CH--CH$_2$--Ph $\xrightarrow{\text{NBS}}$
\qn{4} The number of Cannizaro reaction(s) that occurs in the given transformation is
\qfig{p2_chem_s2_q04_1}
\qn{5} How many resonating structures of carbocation `M' as shown below are possible (excluding given resonating structure) :
\qfig{p2_chem_s2_q05_1}
\qn{6} Find the number of bidentate ligand from the abbreviated ligands given below.\par
dien, bn, ox$^{2-}$, trien, C$_2$O$_2$S$_2^{\,2-}$
\qn{7} How much combinations from the following is/are correct :\par
\noindent{\footnotesize\setlength{\tabcolsep}{2pt}\begin{tabular}{|L{.6cm}|L{.8cm}|L{1.5cm}|L{1.7cm}|L{2.1cm}|}\hline
 & Metal & ore & Chemical Formula & Reduction process\\\hline
(i) & Al & Bauxite & Al$_2$O$_3\cdot$2H$_2$O & Hall's process\\\hline
(ii) & Cu & copper pyrite & CuFeS$_2$ & Bessemerisation\\\hline
(iii) & Fe & Haemetite & Fe$_2$O$_3$ & Smelting\\\hline
(iv) & Pb & Anglesite & PbS & Liquation\\\hline
\end{tabular}}
\qn{8} Find number of reagents which reduces Cr$_2$O$_7^{2-}$ to Cr$^{3+}$ in acidic aqueous medium.\par
SO$_3^{2-}$, CO$_3^{2-}$, H$_2$S, NO$_3^-$, CH$_3$CH$_2$OH, FeSO$_4$, I$^-$, F$^-$, H$^+$/H$_2$O$_2$ $\Delta$
\qn{9} Find the number of acidic radical(s) which can form colourless gas when treated with dilute H$_2$SO$_4$ ?\par
CO$_3^{2-}$, S$^{2-}$, Cl$^-$, SO$_3^{2-}$, HCO$_3^-$, NO$_3^-$, BO$_3^{3-}$
\qn{10} Calculate entropy change (in cal/K) for 1 mole of diatomic ideal gas undergoing for a process A to B as shown in diagram.\par
($R=2$ cal/mole-K ; $\ln2=0.7$)
\qfig{p2_chem_s2_q10_1}
\raggedcolumns
\partsec{3}{MATHEMATICS}
\sect{SECTION-I}
\qn{1} Consider $f(x)=\begin{cases}x^2-2, & x<0\\ x, & 0\le x<2\\ (x-2)^2, & 2\le x<4\\ x-4, & x\ge4\end{cases}$.\par
If the equation $f(x)=k$ has four roots, then $k$ can be equal to :
\begin{qoptions}
\opt{A}{$-2$}
\opt{B}{0}
\opt{C}{2}
\opt{D}{No value of $k$}
\end{qoptions}
\qn{2} Let $f(x)=\begin{cases}\dfrac{1}{|x-1|}, & -\infty<x<0\\ |x|^2+1, & 0\le x<1\\ 2x, & 1\le x<2\\ \dfrac{x^2-1}{x-1}, & \text{otherwise}\end{cases}$, then $f(x)$ is discontinuous in $(-\infty,\infty)$ at
\begin{qoptions}
\opt{A}{two points}
\opt{B}{three points}
\opt{C}{one point}
\opt{D}{no point}
\end{qoptions}
\qn{3} If $y=x+e^x$, then at $x=1$, $\dfrac{d^2x}{dy^2}$ is equal to
\begin{qoptions}
\opt{A}{$e$}
\opt{B}{$\dfrac{-e}{(1+e)^3}$}
\opt{C}{$\dfrac{-e}{(1+e)}$}
\opt{D}{$\dfrac{-e}{(1+e)^2}$}
\end{qoptions}
\qn{4} If $f'(x)$ is non-decreasing function for which $f'(0)=0$, $f(0)=5$ and $f''(x)-f'(x)\le0$ $x\in R$, then the area bounded by the $y=f(x)$ and the lines $x=\dfrac12$ and $x=\dfrac{23}{2}$ does not lies in the interval :
\begin{qoptions}
\opt{A}{[50, 55]}
\opt{B}{[50, 60]}
\opt{C}{[51, 54]}
\opt{D}{[40, 75]}
\end{qoptions}
\qn{5} Let $f(x)=\begin{cases}3x+2\sin^{-1}x+p, & x\le0\\ x^4-x^3+x^2-qx+r, & x>0\end{cases}$.\par
If $f(x)$ is differentiable at $x=0$, $(pr\neq0)$, then
\begin{qoptions}
\opt{A}{$(p-r)^2+(q+5)^2=0$}
\opt{B}{$(p+r)^2+(q+5)^2=0$}
\opt{C}{$(p-r)^2+(q-5)^2=0$}
\opt{D}{$(p+r)^2+(q-5)^2=0$}
\end{qoptions}
\qn{6} The interval on which the function $f(x)=\log_{1/3}(x^2-5x+6)$ is monotonic increasing is
\begin{qoptions}
\opt{A}{$\left(2,\dfrac52\right)\cup(3,\infty)$}
\opt{B}{$(3,\infty)$}
\opt{C}{$(-\infty,2)$}
\opt{D}{$(-\infty,2)\cup\left(\dfrac52,3\right)$}
\end{qoptions}
\qn{7} Consider $f(x)=\begin{cases}x(x+3), & -1\le x<0\\ -\sin x, & 0\le x<\dfrac\pi2\\ -(1+\cos x), & \dfrac\pi2\le x\le\pi\end{cases}$.\par
The sum of maximum and minimum values of $f(x)$ is
\begin{qoptions}
\opt{A}{$-2$}
\opt{B}{$-1$}
\opt{C}{0}
\opt{D}{2}
\end{qoptions}
\qn{8} The equation of normal to the curve $x+y=x^y$, where it cuts x-axis is
\begin{qoptions}
\opt{A}{$y=x+1$}
\opt{B}{$y=-x+1$}
\opt{C}{$y=x-1$}
\opt{D}{$y=-x-1$}
\end{qoptions}
\qn{9} $\displaystyle\int\dfrac{x^4+x^3-1}{x^2\sqrt{x^4+1+2x^3}}\,dx$ is equal to (where $C$ is constant of integration)
\begin{qoptions}
\opt{A}{$\dfrac{\sqrt{x^4+2x^3+1}}{x^{1/2}}+C$}
\opt{B}{$\dfrac{\sqrt{x^4+2x^3+1}}{x}+C$}
\opt{C}{$\dfrac{\sqrt{x^4+2x^3+1}}{x^{3/2}}+C$}
\opt{D}{$\dfrac{\sqrt{x^4+2x^3+1}}{x^3}+C$}
\end{qoptions}
\qn{10} $\displaystyle\int_{\pi/4}^{\pi/2}\left[|\sin x|+|\cos x|+[\operatorname{cosec}x]\right]dx$ is equal to\par
(where $[.]$ denotes greatest integer function)
\begin{qoptions}
\opt{A}{$\dfrac\pi4$}
\opt{B}{$\dfrac\pi2$}
\opt{C}{$\pi$}
\opt{D}{$\dfrac{3\pi}{2}$}
\end{qoptions}
\qn{11} The value of\par\noindent\hfill$\cos\!\bigl(\sin^{-1}\!\bigl(\tan\bigl(\cos^{-1}\bigl(\sin\bigl(\tan^{-1}\tfrac43\bigr)\bigr)\bigr)\bigr)\bigr)$\hfill\null\par\noindent is equal to :
\begin{qoptions}
\opt{A}{$\dfrac35$}
\opt{B}{$\dfrac34$}
\opt{C}{$\dfrac{\sqrt7}{4}$}
\opt{D}{$\dfrac45$}
\end{qoptions}
\qn{12} Let $f(x)=\max.\{x^2-2|x||,\ |x|\}$, then number of points where $f(x)$ is non derivable, is :
\begin{qoptions}
\opt{A}{2}
\opt{B}{3}
\opt{C}{4}
\opt{D}{5}
\end{qoptions}
\qn{13} The tangent of acute angle between the curves $y=|x^2-1|$ and $y=\sqrt{7-x^2}$ at their points of intersection is :
\begin{qoptions}
\opt{A}{$\dfrac{5\sqrt3}{2}$}
\opt{B}{$\dfrac{3\sqrt5}{2}$}
\opt{C}{$\dfrac{5\sqrt3}{4}$}
\opt{D}{$\dfrac{3\sqrt5}{4}$}
\end{qoptions}
\qn{14} The area bounded by the curves $f(x)=9x^2-9x+2$, $g(x)=9x^2-18x+8$ and $x=1$ is :
\begin{qoptions}
\opt{A}{$\dfrac13$}
\opt{B}{$\dfrac12$}
\opt{C}{$\dfrac23$}
\opt{D}{$\dfrac34$}
\end{qoptions}
\qn{15} $\displaystyle\int_{-\pi/4}^{3\pi/4}\dfrac{\left(1+\sin\left(x-\frac\pi4\right)\right)\sqrt2}{\sqrt2-\cos|x|+\sin|x|}\,dx$ is equal to
\begin{qoptions}
\opt{A}{$\ln\left(\dfrac{\sqrt2-1}{\sqrt2}\right)+\sqrt2$}
\opt{B}{$\ln\left(\dfrac{\sqrt2-1}{\sqrt2}\right)+\sqrt2+\dfrac{3\pi}{4}$}
\opt{C}{$\ln\left(\dfrac{\sqrt2}{\sqrt2-1}\right)+\sqrt3+\dfrac{3\pi}{4}$}
\opt{D}{$\ln\left(\dfrac{\sqrt2}{\sqrt2-1}\right)+\sqrt3$}
\end{qoptions}
\qn{16} If $a\sin\theta+b\cos\theta=a\operatorname{cosec}\theta-b\sec\theta=1$, the value of $a^2+b^2+b^{4/3}-b^{2/3}=k$ and $\beta+2\displaystyle\int_0^kx^2e^{-x^2}dx=\int_0^ke^{-x^2}dx$, then the value of $\beta+k$ is equal to
\begin{qoptions}
\opt{A}{$\dfrac{1+e}{e}$}
\opt{B}{$\dfrac{1-e}{e}$}
\opt{C}{0}
\opt{D}{$1+e$}
\end{qoptions}
\qn{17} If $\displaystyle\int\dfrac{x^6(x+1)}{\sqrt{5x^{10}+6x^9+x^4}}\,dx=\dfrac{\sqrt{5x^6+6x^5+1}}{k}+C$ then $k=$
\begin{qoptions}
\opt{A}{10}
\opt{B}{15}
\opt{C}{5}
\opt{D}{None}
\end{qoptions}
\qn{18} $f'(x)=f(x)+\displaystyle\int_0^1f(x)\,dx$ and $f(0)=1$, then $\displaystyle\int f(x)\,dx$ is equal to \underline{\hspace{1cm}} (Where $C$ is constant of integration)
\begin{qoptions}
\opt{A}{$\dfrac{2}{3-e}e^x+\left(\dfrac{3-e}{1-e}\right)x+C$}
\opt{B}{$\dfrac{2}{3-e}e^x+\left(\dfrac{1-e}{3-e}\right)x+C$}
\opt{C}{$\dfrac{2}{2-e}e^x+\left(\dfrac{e}{e-2}\right)x+C$}
\opt{D}{$\left(\dfrac{3}{2-e}\right)e^x+\left(\dfrac{1-e}{2-e}\right)x+C$}
\end{qoptions}
\qn{19} Tangents and normals at points $P$ and $Q$ to the parabola $y^2=4x$ intersect at point $T$ and point $R$ (9, 6) respectively. Then find the length of tangent drawn from $(-1,-1)$ to the circle circumscribing the quadrilateral PTQR.
\begin{qoptions}
\opt{A}{4}
\opt{B}{8}
\opt{C}{3}
\opt{D}{1}
\end{qoptions}
\qn{20} The differential equation of all those lines which are at a unit distance from origin is-
\begin{qoptions}
\opt{A}{$\left(y-x\dfrac{dy}{dx}\right)^2=1-\left(\dfrac{dy}{dx}\right)^2$}
\opt{B}{$\left(y+x\dfrac{dy}{dx}\right)^2=1+\dfrac{dy}{dx}$}
\opt{C}{$\left(y-x\dfrac{dy}{dx}\right)^2=1+\left(\dfrac{dy}{dx}\right)^2$}
\opt{D}{$\left(y+x\dfrac{dy}{dx}\right)^2=1-\left(\dfrac{dy}{dx}\right)^2$}
\end{qoptions}
\sect{SECTION-II}
\qn{1} If the eccentricities of hyperbola $xy=c^2$ and $x^2-y^2=a^2$ are $e_1$ and $e_2$, then the value of $200\,(e_1^{\,4}+e_2^{\,4})$ must be \underline{\hspace{1cm}}
\qn{2} If (5, 12) and (24, 7) are the foci of an ellipse passing through the origin and if $e$ is the eccentricity, then the value of $\left[\dfrac1e\right]$, where $[.]$ represents the greatest integer function is \underline{\hspace{1cm}}
\qn{3} $\displaystyle\int\dfrac{x^4+1}{x^6+1}dx=A\tan^{-1}\left(x-\dfrac1x\right)+B\tan^{-1}\left(x^3\right)+C$ then value of $(A+B)$ is
\qn{4} Let $y=g(x)$ be the image of $f(x)=x+\sin x$ about the line $x+y=0$. If the area bounded by $y=g(x)$, x-axis, $x=0$ and $x=2\pi$ is $A$, then $\dfrac{A}{\pi^2}$ is
\qn{5} If $\displaystyle\int\dfrac{1}{\left(1+\sqrt x\right)^8}dx=\dfrac{-2}{a\left(1+\sqrt x\right)^a}+\dfrac{2}{b\left(1+\sqrt x\right)^b}+C$ (where $C$ is constant of integration), then the value of `a' is equal to
\qn{6} The locus of the mid point of the chords of the ellipse $49x^2+16y^2=784$, the tangents at the ends of which intersect on the circle $x^2+y^2=100$ is $\dfrac{(490x^2+160y^2)^2}{784(x^2+y^2)}=k$ then $k$ is equal to
\qn{7} If the tangent to the conic $y-4=x^2$ at (3, 13) touches the circle $x^2+y^2-4x+2y+\lambda=0$ (for some fixed $\lambda$) at a point $(a,b)$, then $|37a+37b|$ is equal to
\qn{8} Let $f(x)$ is an even function such that $f(x)+f(x-3)=x(x-3)+1$ then $2\displaystyle\int_0^3\dfrac{f(x)}{x^2-3x+1}\,dx$ is
\qn{9} If area bounded by the curve $y=\sqrt{\sin[x]+[\sin x]}$, (where $[.]$ G.I.F.) and lines $x=1$ and $x=\dfrac\pi2$ and x-axis is equal to $\sqrt{\sin1}\left(\dfrac\pi a-b\right)$, then $a+b=$
\qn{10} The area bounded by the curves $|x+y|\le1$, $|y-x|\le1$ and $2x^2+2y^2\ge1$ is $\left|a-\dfrac\pi b\right|$ then $a+b$ is equal to
\end{multicols}
\end{document}
